\documentclass{amsart}
\usepackage{amsmath,amssymb,amsthm,hyperref}

\newtheorem{theorem}{Theorem}
\newtheorem{lemma}{Lemma}
\newtheorem{proposition}{Proposition}
\theoremstyle{definition}
\newtheorem{definition}{Definition}
\newtheorem{remark}{Remark}

\DeclareMathOperator{\Aut}{Aut}
\DeclareMathOperator{\Ham}{Ham}
\DeclareMathOperator{\Lie}{Lie}
\DeclareMathOperator{\grad}{grad}
\DeclareMathOperator{\Scal}{Scal}
\DeclareMathOperator{\Ric}{Ric}

\begin{document}

\title{No non-admissible extremal Sasaki metrics in the join
$M_3\star_{1,101}S^3_{3,2}$}
\maketitle

\section{Statement and answer}

Let $\Sigma_2$ be a closed Riemann surface of genus $2$, and let $M_3$ be the
total space of the primitive (Boothby--Wang) circle bundle over $\Sigma_2$
carrying the regular Sasaki structure whose transverse K\"ahler structure is
the hyperbolic metric of constant Gauss curvature $-1$ on $\Sigma_2$. Let
$S^3_{\mathbf w}$, $\mathbf w=(w_1,w_2)=(3,2)$, denote the $3$-sphere with the
weighted Sasaki structure whose transverse holomorphic structure is the
weighted projective line $\mathbb{CP}^1[3,2]$. Form the Sasaki join
\begin{equation}\label{eq:join}
  M \;=\; M_3 \star_{l_1,l_2} S^3_{\mathbf w},\qquad
  (l_1,l_2)=(1,101),
\end{equation}
a compact $5$-dimensional Sasaki manifold, with its distinguished join Sasaki
structure $(\mathcal S,\bar J)=(\xi,\eta,\Phi,g)$. Recall (see \S\ref{sec:bg})
that the \emph{admissible} Sasaki metrics in the Sasaki cone are those produced
by the join (admissible/Calabi-type) ansatz, i.e.\ those whose transverse
K\"ahler quotient is an \emph{admissible} K\"ahler metric on the
projectivization $P\to\Sigma_2$ compatible with the fibre--base splitting.

\begin{theorem}\label{thm:main}
On $M=M_3\star_{1,101}S^3_{3,2}$ \emph{every} extremal Sasaki metric lying in
the Sasaki class $(\mathcal S,\bar J)$ (equivalently, with fixed transverse
holomorphic structure of the join and Reeb field in the Sasaki cone
$\mathfrak t^+$) is admissible. Consequently the answer to the question is
\textbf{No}: there are no extremal Sasaki metrics in $(\mathcal S,\bar J)$ that
fail to be admissible.
\end{theorem}

The proof rests on facts special to this example:
\begin{itemize}
\item[(A)] the genus of the base is $2$, so $\Sigma_2$ carries
\emph{no} non-zero holomorphic vector field and has finite automorphism group;
consequently \emph{every} transversely holomorphic vector field on $M$ is
vertical (tangent to the $\mathbb{CP}^1[3,2]$ fibres) modulo the Reeb line;
\item[(B)] the maximal torus of the transverse automorphism group of
$(\mathcal S,\bar J)$ then has rank $2$ and the Sasaki cone is exactly the
$2$-dimensional join (``$\mathbf w$-Sasaki'') cone;
\item[(C)] by transverse Calabi symmetry an extremal metric may be taken
invariant under this rank-$2$ torus $T$; closedness of the transverse K\"ahler
form forces the base part of any such metric to be \emph{affine} in the fibre
moment coordinate; and the extremal equation, combined with the vanishing of
holomorphic vector fields on $\Sigma_2$, forces the base K\"ahler metric to have
constant scalar curvature, which on a genus-$2$ surface pins it (uniquely) to
the hyperbolic metric, yielding admissibility.
\end{itemize}
We recall the background and then prove Theorem~\ref{thm:main} via a sequence of
lemmas.

\section{Background and conventions}\label{sec:bg}

\subsection{Sasaki structures and transverse K\"ahler geometry}
A Sasaki structure on a compact manifold $M^{2n+1}$ is a quadruple
$(\xi,\eta,\Phi,g)$ where $\xi$ is the Reeb (unit Killing) field, $\eta$ its
contact form, $\Phi$ an endomorphism with $\Phi^2=-\mathrm{Id}+\xi\otimes\eta$,
and $g$ the associated metric. The Reeb field generates a one-dimensional
foliation $\mathcal F_\xi$ whose transverse geometry is K\"ahler: on the normal
bundle $\nu(\mathcal F_\xi)=TM/\langle\xi\rangle$ one has a complex structure
$\bar J$ (induced by $\Phi$), a transverse K\"ahler form $\frac12 d\eta$, and a
transverse metric $g^T$. The \emph{transverse scalar curvature} $s^T$ is the
scalar curvature of $g^T$. A \emph{Sasaki class} is the set of Sasaki metrics
with fixed $(\xi,\bar J)$ and transverse K\"ahler class $[\frac12 d\eta]$,
deformed by transverse $\partial\bar\partial$-potentials.

\subsection{Extremal Sasaki metrics}
Following Boyer--Galicki--Simanca, a Sasaki metric is \emph{extremal} if the
transverse $(1,0)$-gradient $\partial^\#_g s^T$ of its transverse scalar
curvature is a transversely holomorphic vector field; equivalently $g$ is a
critical point of the transverse Calabi energy $\int_M (s^T)^2\,dV_g$ within
its Sasaki class. A standard identity (the transverse analogue of Calabi's)
shows this is equivalent to $\mathcal L_{\bar J\,\grad^T s^T}\,g^T=0$, i.e.\
$\bar J\,\grad^T s^T$ is a transverse Killing field.

\subsection{The Sasaki cone and the maximal torus}
Let $\Aut(\mathcal S,\bar J)$ be the automorphism group of the transverse
holomorphic structure (transverse biholomorphisms commuting with the flow of
$\xi$), with reduced/identity component $\Aut_0$ whose Lie algebra
$\mathfrak h$ consists of the transversely holomorphic vector fields commuting
with $\xi$. Let $T\subset\Aut_0$ be a maximal torus of rank $r$. The
(unreduced) \emph{Sasaki cone} is
\[
  \mathfrak t^+ \;=\; \{\,\xi'\in\Lie(T) : \eta(\xi')>0 \text{ everywhere}\,\},
\]
an open convex cone of dimension $r$; deforming $\xi$ within $\mathfrak t^+$
keeps $\bar J$ fixed (Boyer--Galicki--Simanca). The relevant search space for
extremal metrics in the problem is exactly $(\mathcal S,\bar J)$: Reeb field in
$\mathfrak t^+$ together with transverse K\"ahler potentials.

\subsection{The join and admissible metrics}\label{ss:admissible}
For \eqref{eq:join} the construction (Boyer--T\o nnesen-Friedman) is the quotient
\[
  M_3\star_{l_1,l_2}S^3_{\mathbf w}
  =\bigl(M_3\times S^3_{\mathbf w}\bigr)/S^1_{(l_2,-l_1)},
\]
where $S^1_{(l_2,-l_1)}$ acts by the difference of the two Reeb circle actions
with the indicated weights; the relative-primality $\gcd(l_1,l_2)=\gcd(101,
w_1w_2)=1$ (here $101$ is prime and coprime to $6$) guarantees the quotient is a
smooth compact Sasaki manifold (a lens-space bundle over $\Sigma_2$), with Reeb
field a positive combination of $\xi_{M_3}$ and $\xi_{\mathbf w}$. In the
quasi-regular case the transverse holomorphic quotient
\[
  P:=M/\mathcal F_\xi
\]
is the total space of a holomorphic $\mathbb{CP}^1[3,2]$ (weighted projective
line / log) bundle $\pi\colon P\to\Sigma_2$ with fibre the weighted projective
line $F:=\mathbb{CP}^1[3,2]$. The \emph{admissible} K\"ahler metrics on $P$ are
the Calabi-type metrics
\begin{equation}\label{eq:admissible}
  g^T \;=\; (1+r\,x)\,g_{\Sigma_2} \;\oplus\; \frac{d x^2}{\Theta(x)} \;\oplus\;
  \Theta(x)\,\theta\otimes\theta,\qquad x\in[-1,1],
\end{equation}
with K\"ahler form
\begin{equation}\label{eq:admissibleform}
  \omega^T=(1+r\,x)\,\omega_{\Sigma_2}+dx\wedge\theta ,\qquad d\theta=r\,\omega_{\Sigma_2},
\end{equation}
where $(g_{\Sigma_2},\omega_{\Sigma_2})$ is the fixed hyperbolic (hence
constant-scalar-curvature) K\"ahler metric on $\Sigma_2$, $x$ is the moment
coordinate of the fibrewise $S^1$ action, $\theta$ a connection $1$-form on the
principal part with curvature $r\,\omega_{\Sigma_2}$ ($r$ a fixed real constant
determined by the bundle, with $1+rx>0$ on $[-1,1]$), and $\Theta>0$ on
$(-1,1)$ with $\Theta(\pm1)=0$, $\Theta'(\pm1)=\mp2$, a smooth profile.
Equivalently: the admissible metrics are exactly the metrics that are invariant
under the fibrewise torus, split as base $\oplus$ fibre, have base part a
\emph{fixed constant-scalar-curvature} K\"ahler metric scaled affinely in $x$,
and fibre data depending only on $x$.

\section{Rigidity of the base and verticality}

\begin{lemma}\label{lem:rigidbase}
$H^0(\Sigma_2,T^{1,0}\Sigma_2)=0$: $\Sigma_2$ admits no non-zero holomorphic
vector field, and $\Aut(\Sigma_2)$ is finite. Moreover the constant-scalar-curvature
K\"ahler metric in each K\"ahler class on $\Sigma_2$ is unique up to (the finite
group of) automorphisms and is the hyperbolic metric of that area.
\end{lemma}

\begin{proof}
On a compact Riemann surface of genus $g$ the holomorphic vector fields are
$H^0(\Sigma,T^{1,0}\Sigma)=H^0(\Sigma,K_\Sigma^{-1})$, where $K_\Sigma$ is the
canonical bundle of degree $2g-2$, so $\deg K_\Sigma^{-1}=2-2g$. For $g=2$ this
is $-2<0$. A line bundle of negative degree on a compact Riemann surface
has no non-zero holomorphic section (a non-zero section would have an effective
divisor of degree equal to the bundle degree, which would be $\ge 0$). Hence
$H^0(\Sigma_2,T^{1,0}\Sigma_2)=0$. The identity component
$\Aut_0(\Sigma_2)$ is a complex Lie group with Lie algebra
$H^0(\Sigma_2,T^{1,0}\Sigma_2)=0$, hence trivial; together with Hurwitz's
finiteness of $\Aut(\Sigma_2)$ for $g\ge2$, $\Aut(\Sigma_2)$ is finite.

For the last assertion: on a compact Riemann surface a constant-scalar-curvature
K\"ahler metric is exactly a constant-Gauss-curvature metric. Fix a K\"ahler
class, i.e.\ fix the total area $A>0$. By the uniformization theorem the unique
constant-Gauss-curvature metric in the conformal class of $\Sigma_2$ has Gauss
curvature $K_0<0$ and area $A_0=2\pi|\chi(\Sigma_2)|/|K_0|$; rescaling by
$A/A_0$ produces the unique constant-curvature representative of total area $A$
in that conformal class. Any constant-scalar-curvature K\"ahler metric of area
$A$ is conformal to this one, say $g'=e^{2u}g_0$ with $g_0$ the
constant-curvature metric and $K'=$ const; the Gauss-curvature transformation
law $K'=e^{-2u}(K_0-\Delta_{g_0}u)$ with $K_0,K'$ constant and the area
constraint $\int e^{2u}\,dA_0=A=\int dA_0$ force, by the maximum principle
applied to $u$ at its extrema in $\Delta_{g_0}u=K_0-K'e^{2u}$, that $u\equiv0$;
hence $g'=g_0$. Uniqueness up to $\Aut(\Sigma_2)$ follows.
\end{proof}

\begin{lemma}\label{lem:vertical}
Every holomorphic vector field on the total space $P$ of the fibration
$\pi\colon P\to\Sigma_2$ (with rational fibre $F=\mathbb{CP}^1[3,2]$) is
vertical, i.e.\ tangent to the fibres:
$H^0(P,T_P)=H^0(P,T_{P/\Sigma_2})$. Consequently every transversely holomorphic
vector field on $M$ commuting with $\xi$ is, modulo the Reeb line, vertical.
\end{lemma}

\begin{proof}
Differentiating $\pi$ gives the short exact sequence of holomorphic sheaves
\[
  0 \longrightarrow T_{P/\Sigma_2} \longrightarrow T_P
  \xrightarrow{\ d\pi\ } \pi^*T_{\Sigma_2} \longrightarrow 0 ,
\]
where $T_{P/\Sigma_2}=\ker d\pi$ is the relative (vertical) tangent sheaf. Taking
global sections yields the left-exact sequence
\[
  0 \to H^0(P,T_{P/\Sigma_2}) \to H^0(P,T_P)
  \xrightarrow{\ d\pi\ } H^0(P,\pi^*T_{\Sigma_2}).
\]
By the projection formula,
$H^0(P,\pi^*T_{\Sigma_2})=H^0\!\bigl(\Sigma_2,T_{\Sigma_2}\otimes
\pi_*\mathcal O_P\bigr)$. Since the fibre $F$ is a compact connected rational
curve (weighted projective line), $H^0(F,\mathcal O_F)=\mathbb C$ and
$\pi_*\mathcal O_P=\mathcal O_{\Sigma_2}$ (the fibres are connected with only
constants as global functions, and the bundle is locally trivial in the
analytic/orbifold sense). Hence
\[
  H^0(P,\pi^*T_{\Sigma_2})=H^0(\Sigma_2,T_{\Sigma_2})=0
\]
by Lemma~\ref{lem:rigidbase}. Therefore the map
$d\pi\colon H^0(P,T_P)\to H^0(P,\pi^*T_{\Sigma_2})$ is the zero map, and
$H^0(P,T_P)=H^0(P,T_{P/\Sigma_2})$: every holomorphic vector field on $P$ is
vertical.

For the Sasaki statement: a transversely holomorphic vector field on $M$
commuting with $\xi$ descends (in the quasi-regular case to $P$ exactly, and in
general by approximating the Reeb field by quasi-regular ones in the same Sasaki
cone and passing to the limit, the descended holomorphic field on $P$ varying
continuously) to a holomorphic vector field on $P=M/\mathcal F_\xi$, modulo the
Reeb generator $\xi$ itself. By the previous paragraph this descended field is
vertical, which is the assertion.
\end{proof}

\begin{lemma}\label{lem:cone}
For $M=M_3\star_{1,101}S^3_{3,2}$, the maximal torus of $\Aut_0(\mathcal S,\bar
J)$ has rank $2$, and the Sasaki cone $\mathfrak t^+$ is the $2$-dimensional
join cone, spanned by $\xi_{M_3}$ and $\xi_{\mathbf w}$.
\end{lemma}

\begin{proof}
By Lemma~\ref{lem:vertical}, $\mathfrak h=\Lie\Aut_0(\mathcal S,\bar J)$ is
spanned by the Reeb direction $\xi$ together with the vertical (fibrewise)
holomorphic vector fields of $P\to\Sigma_2$. The fibre $F=\mathbb{CP}^1[3,2]$ is
a weighted projective line with distinct weights $w_1=3\neq w_2=2$; its
connected holomorphic automorphism group $\Aut_0(F)$ is the image of the
$\mathbb C^\times$ acting with weights $(3,2)$, i.e.\ a one-dimensional torus
$\mathbb C^\times$ fixing the two orbifold points $[1:0]$ and $[0:1]$ (the
``triangular'' unipotent flows present for the ordinary $\mathbb{CP}^1$ are
absent here because the two weights are distinct: a global vector field on
$\mathbb{CP}^1[3,2]$ corresponds to a weighted-homogeneous vector field on
$\mathbb C^2\setminus\{0\}$ of weighted degree $0$ for the $(3,2)$-action, and
the only such fields are scalar multiples of the Euler field; so no
$\mathfrak{sl}_2$ appears). Hence the space of vertical holomorphic vector
fields is the rank-one relative one whose maximal compact subgroup is a
\emph{single} fibrewise circle $S^1_{\mathrm{fib}}$; there is no further compact
toric direction because any additional commuting field would have to be
non-vertical, contradicting Lemma~\ref{lem:vertical}, or would require an
$\mathfrak{sl}_2$ in the fibre, which $\mathbb{CP}^1[3,2]$ does not have.

Therefore a maximal torus $T$ of $\Aut_0(\mathcal S,\bar J)$ is generated by the
Reeb circle and $S^1_{\mathrm{fib}}$, of rank
$r = 1\ (\text{Reeb}) + 1\ (\text{fibre}) = 2$.
All maximal tori are conjugate, and the two generators are, up to the join
identification, $\xi_{M_3}$ and $\xi_{\mathbf w}$; hence
$\mathfrak t^+=\{\xi'\in\Lie(T):\eta(\xi')>0\}$ is the $2$-dimensional join
($\mathbf w$-Sasaki) cone.
\end{proof}

\section{Calabi symmetry and the structure of invariant metrics}

\begin{lemma}[Transverse Calabi symmetry]\label{lem:average}
Let $g$ be an extremal Sasaki metric in $(\mathcal S,\bar J)$ with Reeb field
$\xi\in\mathfrak t^+$. Then after a transverse biholomorphism the transverse
metric $g^T$ is invariant under the maximal torus $T$ of Lemma~\ref{lem:cone},
and the extremal field $V:=\partial^\#_g s^T$ lies in $\Lie(T)\otimes\mathbb C$.
Moreover $V$ is \emph{vertical}, and consequently $s^T$ is a function of the
fibre moment coordinate $x$ alone.
\end{lemma}

\begin{proof}
\emph{Step 1: the extremal field generates transverse isometries.} By
extremality $V=\partial^\#_g s^T$ is transversely holomorphic, and the
identity recalled in \S\ref{sec:bg} gives $\mathcal L_{\bar
J\grad^T s^T}g^T=0$, i.e.\ $X:=\bar J\,\grad^T s^T=\mathrm{Im}\,V$ is a
transverse Killing field. Its flow is a one-parameter group of transverse
isometries; its closure $T_V$ is a torus inside the (compact) isometry group
$\mathrm{Isom}(g^T)$ of the transverse metric.

\emph{Step 2: maximal compact symmetry (transverse Calabi theorem).} By the
transverse Sasaki analogue of Calabi's theorem on extremal K\"ahler metrics
(Boyer--Galicki--Simanca, ``Canonical Sasakian metrics'', Comm.\ Math.\ Phys.\
279 (2008); transcribing Calabi, ``Extremal K\"ahler metrics II'', 1985, to the
transverse K\"ahler foliation), the identity component of $\mathrm{Isom}(g^T)$
for an extremal transverse K\"ahler metric is a \emph{maximal compact subgroup}
$K$ of $\Aut_0(\mathcal S,\bar J)$, and the extremal field lies in the centre of
$\Lie K$; in particular $V$ lies in the complexified Lie algebra of a maximal
torus $T'\subset K$. The hypotheses (compactness of $M$; $g$ extremal in a fixed
Sasaki class) hold here.

\emph{Step 3: identification of the torus.} By Lemmas~\ref{lem:vertical}
and~\ref{lem:cone} every maximal torus of $\Aut_0(\mathcal S,\bar J)$ has rank
$2$ (Reeb $+$ fibre) and all such tori are conjugate. Conjugating by a
transverse biholomorphism we arrange $T'=T$; then $g^T$ is $T$-invariant and
$V\in\Lie(T)\otimes\mathbb C$.

\emph{Step 4: verticality and $s^T=s^T(x)$.} The complexified Lie algebra
$\Lie(T)\otimes\mathbb C$ is spanned by the Reeb direction $\xi$ (which acts
trivially on the transverse geometry) and the fibre generator. Hence, as a
transverse holomorphic vector field, $V$ is a (complex) multiple of the
fibrewise generator of $S^1_{\mathrm{fib}}$, i.e.\ \emph{vertical}: $d\pi(V)=0$.
By definition $V=\partial^\#_g s^T$ is the metric dual of $\partial s^T$, so for
every transverse vector $W$,
\[
  g^T(V,\bar W)=\langle \partial s^T, \bar W\rangle = \tfrac12\bigl(W s^T - i\,(\bar J W)s^T\bigr).
\]
Since $S^1_{\mathrm{fib}}$ acts by isometries fixing the base, the bundle
splitting into the vertical (fibre) and horizontal distributions is
$g^T$-orthogonal (the horizontal distribution being the $g^T$-orthogonal
complement of the fibre tangent space). Choosing $W$ horizontal, orthogonality
gives $g^T(V,\bar W)=0$, hence $W s^T=(\bar J W)s^T=0$ for all horizontal $W$;
i.e.\ $ds^T$ annihilates the horizontal distribution and $s^T$ is constant along
the base. Being also $S^1_{\mathrm{fib}}$-invariant, $s^T$ is a function of the
single fibre moment coordinate $x\in[-1,1]$ alone.
\end{proof}

\begin{lemma}[Affine base from closedness]\label{lem:affine}
Let $g^T$ be a $T$-invariant transverse K\"ahler metric on $P\to\Sigma_2$ with
Reeb field in $\Lie(T)$. Then on the open dense set $P^\circ=\{-1<x<1\}$ its
K\"ahler form is
\begin{equation}\label{eq:Tinv-form}
  \omega^T \;=\; \omega_x \;+\; dx\wedge\theta ,
\end{equation}
where $\theta$ is a connection $1$-form for $S^1_{\mathrm{fib}}$
($\theta(\partial_t)=1$ on the fibre circle, $\theta|_{\mathcal H}=0$ on the
$g^T$-horizontal distribution $\mathcal H$) with $d\theta=\rho$
($\rho=\pi^*\bar\rho$ for a fixed closed $2$-form $\bar\rho$ on $\Sigma_2$
representing the curvature of the bundle), $x$ the $S^1_{\mathrm{fib}}$-moment
map, and $\omega_x$ a smooth family of K\"ahler forms on $\Sigma_2$ satisfying
\begin{equation}\label{eq:affine}
  \omega_x \;=\; \omega_{0} \;+\; x\,\rho ,
\end{equation}
i.e.\ \emph{affine in $x$ with the fixed slope $\rho$}, $\omega_0$ a single
K\"ahler form on $\Sigma_2$.
\end{lemma}

\begin{proof}
Let $x$ be the moment map of $S^1_{\mathrm{fib}}$ for $\omega^T$, normalised to
take values in $[-1,1]$ (its image is the closed interval, with the two
endpoints the fixed sections, by compactness of the fibre and the standard
normalisation of the weighted-$\mathbb{CP}^1$ moment map). Let $\theta$ be as
stated and $\rho:=d\theta$. Because $S^1_{\mathrm{fib}}$ acts by holomorphic
isometries fixing the base, $\mathcal H=\ker\theta\cap\ker dx$ projects
isomorphically to $T\Sigma_2$ at each point and is $g^T$-orthogonal to the
fibre. The defining relation of the moment map is
$\iota_{\partial_t}\omega^T=dx$, which together with $T$-invariance forces the
mixed $\partial_t$-components of $\omega^T$ to be exactly $dx\wedge\theta$; the
remaining (horizontal) part is a $T$-invariant $2$-form pulled back from
$\Sigma_2$ for each fixed $x$, denoted $\omega_x$. Positivity of $g^T$ makes
$\omega_x$ a K\"ahler form on $\Sigma_2$. This gives \eqref{eq:Tinv-form}.

Impose $d\omega^T=0$. With $d(dx\wedge\theta)=-dx\wedge d\theta=-dx\wedge\rho$
and $d=d_{\Sigma}+dx\wedge\partial_x$ on $S^1_{\mathrm{fib}}$-invariant forms
($\rho$ pulled back from $\Sigma_2$, so $\partial_x\rho=0$),
\[
  0=d\omega^T = d_\Sigma\omega_x + dx\wedge\partial_x\omega_x - dx\wedge\rho .
\]
The term $d_\Sigma\omega_x$ is a $3$-form on the $2$-real-dimensional base
$\Sigma_2$, hence $\equiv0$; matching the $dx\wedge(\cdot)$ components gives
$\partial_x\omega_x=\rho$, a closed $2$-form independent of $x$. Integrating in
$x$ from $0$ gives $\omega_x=\omega_0+x\,\rho$, which is \eqref{eq:affine}. Since
$[\rho]=[d\theta]$ is the Chern class of the $S^1_{\mathrm{fib}}$-bundle, $\rho$
is the fixed topological curvature form.
\end{proof}

\section{Extremality forces the base to be CSC}

We first record the transverse scalar curvature of the structurally constrained
metrics produced by Lemmas~\ref{lem:average}--\ref{lem:affine}.

\begin{lemma}[Scalar curvature of an invariant affine-base metric]\label{lem:scal}
Let $g^T$ be as in \eqref{eq:Tinv-form}--\eqref{eq:affine}, written in the
orthogonal block form
\begin{equation}\label{eq:block}
  g^T = g_x \ \oplus\ \frac{1}{\Theta(x)}\,dx^2 \ \oplus\ \Theta(x)\,\theta\otimes\theta ,
\end{equation}
where $g_x$ is the base K\"ahler metric with K\"ahler form $\omega_x$ and
$\Theta(x)=g^T(\partial_t,\partial_t)>0$ on $(-1,1)$. Fix the hyperbolic
reference $(g_{\Sigma_2},\omega_{\Sigma_2})$ and write the positive density
\[
  F(b,x):=\frac{\omega_x}{\omega_{\Sigma_2}}(b)=a(b)+x\,c(b),\qquad
  a:=\frac{\omega_0}{\omega_{\Sigma_2}},\quad c:=\frac{\rho}{\omega_{\Sigma_2}} .
\]
Then the transverse scalar curvature is
\begin{equation}\label{eq:scalformula}
  s^T(b,x)\;=\;\Scal(g_x)(b)\;-\;\frac{\bigl(F(b,x)\,\Theta(x)\bigr)''}{F(b,x)},
\end{equation}
where $\,'=d/dx$, and the base scalar curvature satisfies the conformal identity
\begin{equation}\label{eq:confscal}
  \Scal(g_x)(b)\;=\;\frac{1}{F(b,x)}\Bigl(\Scal(g_{\Sigma_2})-\Delta_{g_{\Sigma_2}}\log F(b,x)\Bigr),
\end{equation}
with $\Scal(g_{\Sigma_2})=-2$ (Gauss curvature $-1$, real dimension $2$) and
$\Delta_{g_{\Sigma_2}}$ the Laplace--Beltrami operator (so $\Delta f\le0$ at an
interior maximum of $f$).
\end{lemma}

\begin{proof}
Formula \eqref{eq:scalformula} is the scalar curvature of the Calabi
(``momentum'') construction over a K\"ahler base, specialised to a base of
complex dimension one. It is obtained as follows. The block form
\eqref{eq:block} is the Riemannian submersion metric of the principal
$S^1_{\mathrm{fib}}$-bundle over $\Sigma_2\times(-1,1)$ with fibre coordinate
$t$, $\iota_{\partial_t}\omega^T=dx$, fibre length squared $\Theta(x)$, and
curvature $\rho$; this is exactly the data of the generalized Calabi ansatz of
Apostolov--Calderbank--Gauduchon--T\o nnesen-Friedman (``Hamiltonian $2$-forms
in K\"ahler geometry, II'', J.\ Differential Geom.\ 68 (2004); see also their
``Extremal K\"ahler metrics on projective bundles over a curve'', Adv.\ Math.\
227 (2011)). The general scalar-curvature formula there, for a base K\"ahler
product $\prod_a(B_a,\pm g_a)$ with momentum profile $\Theta$ and characteristic
polynomial $p(x)=\prod_a(1+ r_a x)^{d_a}$, is
$\,\Scal=\sum_a \frac{\Scal(g_a)}{1+r_a x}-\frac1{p(x)}\bigl(p(x)\Theta(x)\bigr)''$.
In our case there is a single base factor $(\Sigma_2,g_{\Sigma_2})$ and the
``$1+r_ax$'' factor is replaced by the genuine pointwise density
$F(b,x)=\omega_x/\omega_{\Sigma_2}$ (which reduces to $1+rx$ exactly in the
admissible case $\omega_x=(1+rx)\omega_{\Sigma_2}$); the first summand becomes
the actual scalar curvature $\Scal(g_x)$ of the rescaled base metric and the
fibre term becomes $(F\Theta)''/F$. This is \eqref{eq:scalformula}.

Identity \eqref{eq:confscal} is the conformal change of Gauss curvature on a
surface: if $g_x=F\,g_{\Sigma_2}=e^{2\varphi}g_{\Sigma_2}$ with
$\varphi=\tfrac12\log F$, then the Gauss curvatures satisfy
$K_x=e^{-2\varphi}\bigl(K_{\Sigma_2}-\Delta_{g_{\Sigma_2}}\varphi\bigr)$;
multiplying by $2$ ($\Scal=2K$) and using $2\varphi=\log F$ gives
\eqref{eq:confscal}.
\end{proof}

\begin{lemma}\label{lem:baseCSC}
Let $g^T$ be a $T$-invariant \emph{extremal} transverse K\"ahler metric on
$P\to\Sigma_2$ as in Lemmas~\ref{lem:average}--\ref{lem:scal}. Then the density
$u:=c/a=\rho/\omega_0$ is constant on $\Sigma_2$; equivalently $\rho$ is a
constant multiple of $\omega_0$ (and of every $\omega_x$). Consequently the base
metric $g_x$ has constant scalar curvature, $\omega_0$ and $\rho$ are constant
multiples of the hyperbolic form $\omega_{\Sigma_2}$, and $g^T$ is admissible
\eqref{eq:admissible}.
\end{lemma}

\begin{proof}
By Lemma~\ref{lem:average}, extremality (after the Calabi symmetrisation) is
equivalent to $s^T$ being a function of $x$ \emph{alone}: $s^T(b,x)=\sigma(x)$
for some smooth $\sigma$, for all $b\in\Sigma_2$ and all $x\in(-1,1)$. Insert
\eqref{eq:scalformula} and multiply by $F=a+xc$:
\begin{equation}\label{eq:master}
  F\,\sigma(x)\;=\;F\,\Scal(g_x)\;-\;(F\Theta)''
  \;=\;\bigl(\Scal(g_{\Sigma_2})-\Delta_{g_{\Sigma_2}}\log F\bigr)-(F\Theta)'' ,
\end{equation}
using \eqref{eq:confscal}. Write $\Scal(g_{\Sigma_2})=:k$ (a constant, $=-2$)
and expand $(F\Theta)''=(a+xc)\Theta''+2c\Theta'$ (recall $a,c$ are functions of
$b$ only, $\Theta$ of $x$ only). Then \eqref{eq:master} reads, for all $b,x$,
\begin{equation}\label{eq:master2}
  \Delta_{g_{\Sigma_2}}\log(a+xc)
  \;=\; k - a\bigl(\sigma(x)+\Theta''(x)\bigr)
        \;-\; c\bigl(x\sigma(x)+x\Theta''(x)+2\Theta'(x)\bigr).
\end{equation}
The right-hand side is, for each fixed $x$, an affine combination of the two
fixed base functions $a$ and $c$ with $b$-independent coefficients.

\emph{Power-series matching in $x$.} Rearrange \eqref{eq:master2}. Since
$\log(a+xc)=\log a+\log(1+xu)$ with $u:=c/a$, and $\log a$ is $x$-independent,
\begin{equation}\label{eq:masterU}
  \Delta_{g_{\Sigma_2}}\log(1+xu)
  \;=\; G \;-\; a\,S(x)\;-\;a\,u\,T(x),
\end{equation}
where $G:=k-\Delta_{g_{\Sigma_2}}\log a$ is a function of $b$ alone,
$S(x):=\sigma(x)+\Theta''(x)$ and $T(x):=x\sigma(x)+x\Theta''(x)+2\Theta'(x)$ are
$b$-independent smooth functions of $x$, and we used $c=a\,u$. Both sides of
\eqref{eq:masterU} are smooth in $x$ on $(-1,1)$, so \eqref{eq:masterU} holds
together with all its $x$-derivatives at $x=0$. Using the convergent expansion
$\log(1+xu)=\sum_{n\ge1}\frac{(-1)^{n+1}}{n}(xu)^n$ for small $x$, and that
$\Delta_{g_{\Sigma_2}}$ commutes with $\partial_x$,
\[
  \Delta_{g_{\Sigma_2}}\log(1+xu)=\sum_{n\ge1}\frac{(-1)^{n+1}}{n}\,x^n\,
  \Delta_{g_{\Sigma_2}}(u^n).
\]
Matching the coefficient of $x^1$ in \eqref{eq:masterU} (equivalently,
$\partial_x|_{x=0}$) gives, with $A_1:=S'(0)$ and $B_1:=T'(0)$ constants,
\begin{equation}\label{eq:E1}
  \Delta_{g_{\Sigma_2}}u \;=\; -a\,(A_1+B_1\,u),
\end{equation}
since the right-hand side contributes $-a\,S'(0)-a\,u\,T'(0)$ at first order in
$x$ ($G$ being constant in $x$). Matching the coefficient of $x^2$
(equivalently, $\tfrac12\partial_x^2|_{x=0}$) gives, with $A_2:=S''(0)$ and
$B_2:=T''(0)$,
\[
  -\tfrac12\,\Delta_{g_{\Sigma_2}}(u^2)\;=\;-\tfrac12\,a\,(A_2+B_2\,u),
  \qquad\text{i.e.}\qquad \Delta_{g_{\Sigma_2}}(u^2)=a\,(A_2+B_2\,u).
\]
Substituting $\Delta_{g_{\Sigma_2}}(u^2)=2u\,\Delta_{g_{\Sigma_2}}u+2|\nabla
u|^2$ and \eqref{eq:E1}, and solving for $|\nabla u|^2$ (all norms with respect
to $g_{\Sigma_2}$), we obtain
\begin{equation}\label{eq:E2}
  |\nabla u|^2 \;=\; a\,Q_2(u),\qquad
  Q_2(u)=B_1\,u^2+\bigl(A_1+\tfrac12 B_2\bigr)u+\tfrac12 A_2 .
\end{equation}
Crucially the leading coefficient of $Q_2$ is the \emph{same} constant $B_1$ as
in \eqref{eq:E1}. Equations \eqref{eq:E1}--\eqref{eq:E2} are the only inputs we
need.

\emph{Maximum principle.} Suppose, for contradiction, that $u$ is not constant.
Let $U=\max_{\Sigma_2}u$ attained at $p$ and $L=\min_{\Sigma_2}u<U$ attained at
$q$. At $p$, $\nabla u=0$, so \eqref{eq:E2} gives $a(p)Q_2(U)=0$, hence
$Q_2(U)=0$ (as $a>0$); similarly $Q_2(L)=0$. Since $U\ne L$, the quadratic $Q_2$
has the two distinct roots $L,U$, so $Q_2(s)=B_1(s-U)(s-L)$. For interior values
$s=u(b)\in(L,U)$ we have $(s-U)(s-L)<0$; since $|\nabla u|^2\ge0$ and $a>0$,
\eqref{eq:E2} forces $B_1\le0$. If $B_1=0$ then $Q_2\equiv0$, so $|\nabla
u|^2\equiv0$ and $u$ is constant, a contradiction; hence $B_1<0$.

Now use \eqref{eq:E1}: with $u_\ast:=-A_1/B_1$,
\[
  \Delta_{g_{\Sigma_2}}u \;=\; -a\,B_1\,(u-u_\ast)\;=\;(-aB_1)\,(u-u_\ast),
  \qquad -aB_1>0 .
\]
At the maximum $p$: $\Delta_{g_{\Sigma_2}}u(p)\le0$, so $(-a(p)B_1)(U-u_\ast)\le0$,
giving $U\le u_\ast$. At the minimum $q$: $\Delta_{g_{\Sigma_2}}u(q)\ge0$, so
$(-a(q)B_1)(L-u_\ast)\ge0$, giving $L\ge u_\ast$. Combining, $U\le u_\ast\le L$,
contradicting $L<U$. Therefore $u$ is constant.

\emph{Conclusion.} With $u=c/a$ constant, $\rho=u\,\omega_0$, so $\rho$ is a
constant multiple of $\omega_0$, and $\omega_x=(1+xu)\,\omega_0$ is a constant
multiple of $\omega_0$ for each $x$; in particular $F(b,x)=(1+xu)\,a(b)$ and
$\log F=\log(1+xu)+\log a$, whence $\Delta_{g_{\Sigma_2}}\log F
=\Delta_{g_{\Sigma_2}}\log a$ is $x$-independent. Plugging back into
\eqref{eq:confscal}, $\Scal(g_x)=\frac{1}{(1+xu)a}\,(k-\Delta_{g_{\Sigma_2}}\log
a)$. But $s^T=\sigma(x)$ depends on $x$ only and, by \eqref{eq:scalformula} with
$F=(1+xu)a$, the fibre term $(F\Theta)''/F=(a(1+xu)\Theta)''/(a(1+xu))$ is also a
function of $x$ only times $1/a\cdot a=1$, i.e.\ independent of $b$ once $u$ is
constant; hence $\Scal(g_x)$ must be independent of $b$, i.e.\
$k-\Delta_{g_{\Sigma_2}}\log a\equiv\text{const}$, so $\Delta_{g_{\Sigma_2}}\log
a$ is constant. On the compact $\Sigma_2$, $\int_{\Sigma_2}\Delta_{g_{\Sigma_2}}
\log a\,\,dV_{g_{\Sigma_2}}=0$, so the constant is $0$ and $\Delta_{g_{\Sigma_2}}
\log a\equiv0$; by the maximum principle (a function with vanishing Laplacian on
a closed manifold is constant) $\log a$ is constant, i.e.\ $a$ is a positive
constant. Thus $\omega_0=a\,\omega_{\Sigma_2}$ and
$\rho=u\,a\,\omega_{\Sigma_2}$ are constant multiples of the hyperbolic form;
$g_x=(1+ux)a\,g_{\Sigma_2}$ has constant scalar curvature for each $x$, and by
Lemma~\ref{lem:rigidbase} this is the (unique) hyperbolic CSC metric.

Renormalising so that $a=1$ (absorbing the constant scale into
$\omega_{\Sigma_2}$) and setting $r:=u$, the metric \eqref{eq:block} becomes
\[
  \omega^T=(1+r\,x)\,\omega_{\Sigma_2}+dx\wedge\theta,\qquad d\theta=\rho=r\,\omega_{\Sigma_2},
\]
which is exactly the admissible form \eqref{eq:admissibleform}; the fibre profile
$\Theta(x)=g^T(\partial_t,\partial_t)$ depends on $x$ only and satisfies the
boundary conditions $\Theta(\pm1)=0,\ \Theta'(\pm1)=\mp2$ forced by smooth
collapse of the fibre at the two end sections $x=\pm1$. Hence $g^T$ is
admissible.
\end{proof}

\section{Proof of Theorem~\ref{thm:main}}

Let $g$ be any extremal Sasaki metric in the class $(\mathcal S,\bar J)$, with
Reeb field $\xi\in\mathfrak t^+$.

By Lemma~\ref{lem:average}, after a transverse biholomorphism the transverse
K\"ahler metric $g^T$ is invariant under the maximal torus $T$ of
Lemma~\ref{lem:cone}, its Reeb field lies in $\Lie(T)$, the extremal field
$V=\partial^\#_g s^T$ is vertical, and $s^T=s^T(x)$ depends only on the fibre
moment coordinate.

By Lemma~\ref{lem:affine}, $T$-invariance together with closedness of the
transverse K\"ahler form forces the base part of $g^T$ to be affine in $x$ with
the fixed topological slope $\rho$: $\omega_x=\omega_0+x\rho$, and the metric
takes the orthogonal block form \eqref{eq:block} with scalar curvature given by
Lemma~\ref{lem:scal}.

By Lemma~\ref{lem:baseCSC}, the requirement $s^T=s^T(x)$ (extremality) then
forces $\omega_0$ and $\rho$ to be constant multiples of the (unique, by
Lemma~\ref{lem:rigidbase}) hyperbolic CSC form $\omega_{\Sigma_2}$, and the
metric to be of the admissible form
\eqref{eq:admissibleform}--\eqref{eq:admissible}. Hence $g$ is admissible.

Since $g$ was an arbitrary extremal Sasaki metric in $(\mathcal S,\bar J)$,
every extremal Sasaki metric in this class is admissible. Equivalently,
\emph{there are no} extremal Sasaki metrics in $(\mathcal S,\bar J)$ that fail
to be admissible. This proves Theorem~\ref{thm:main}. \qed

\begin{remark}[On the role of genus $\ge2$]
The genus-$2$ hypothesis enters \emph{twice} and is essential to the answer
``No'': (i) in Lemma~\ref{lem:rigidbase}, to kill all holomorphic vector fields
on $\Sigma_2$, which is what makes every transverse holomorphic field vertical
(Lemma~\ref{lem:vertical}), the Sasaki cone exactly rank $2$
(Lemma~\ref{lem:cone}), and the extremal field vertical so that $s^T=s^T(x)$
(Lemma~\ref{lem:average}); and (ii) in Lemma~\ref{lem:baseCSC}, where the
maximum-principle argument forces the base density to be constant and, via
uniqueness of the hyperbolic metric, pins the base to the admissible CSC metric.
For a base of genus $0$ or $1$ (where holomorphic vector fields exist), the
extremal field need not be vertical and the base part need not be CSC, leaving
room for genuinely non-admissible extremal metrics; the result is therefore
special to higher genus.
\end{remark}

\begin{remark}[On existence]
Theorem~\ref{thm:main} is a uniqueness/rigidity statement: it asserts that the
extremal locus, if non-empty, lies inside the admissible family. Existence of an
extremal (indeed constant-scalar-curvature, when the relevant
Boyer--T\o nnesen-Friedman quasi-polynomial has the appropriate sign/roots)
admissible representative is a separate, ODE-level question for the profile
$\Theta$ within \eqref{eq:admissible}; it does not affect the negative answer to
the posed question, which concerns whether \emph{non-admissible} extremal
metrics can occur.
\end{remark}

\end{document}
